%%%%%%%%%%%%%%%%%%%%%%%%%%%%%%%%%%%%%%%%%%%%%%%%%%%%%%%%%%%%%
% ==================== 第四章：符号说明 ====================
% BZD 数模社制作——国赛 B 题通用模板
% 使用时请根据全文实际采用的符号替换【】中的内容，并删除多余行。
%
% ✓ 自查清单：
%   □ 是否采用三线表，并按集合/索引、变量、参数、指标的逻辑排列？
%   □ 是否只收录全文反复使用的主要符号？
%   □ 是否删除仅出现一次、可在公式后直接解释的局部符号？
%   □ 有量纲变量是否标明单位，无量纲量是否明确写“无量纲”？
%   □ 同一符号是否只表示一个含义？
%   □ 同一变量是否始终使用同一符号、上下标和单位？
%   □ 是否检查了题面、模型、代码和结果表之间的符号与单位一致性？
%%%%%%%%%%%%%%%%%%%%%%%%%%%%%%%%%%%%%%%%%%%%%%%%%%%%%%%%%%%%%

\section{符号说明}

为便于后续模型表达，现将全文反复使用的主要符号汇总如表~\ref{tab:symbols} 所示。仅在局部公式中出现的符号，将在正文首次出现时另行说明。

\begin{table}[htbp]
  \centering
  \caption{主要符号说明}
  \label{tab:symbols}
  \renewcommand{\arraystretch}{1.25}
  \begin{tabular}{@{}
    >{\centering\arraybackslash}p{0.18\textwidth}
    >{\raggedright\arraybackslash}p{0.56\textwidth}
    >{\centering\arraybackslash}p{0.14\textwidth}@{}}
    \toprule
    \textbf{符号} & \centering\textbf{含义}\arraybackslash & \textbf{单位} \\
    \midrule
    【集合或索引符号】 & 【集合、对象编号、阶段编号或时间索引的含义】 & 【无量纲】 \\
    【状态或观测变量】 & 【描述研究对象状态、位置、数量或观测结果的变量含义】 & 【实际单位】 \\
    【决策变量】 & 【模型需要确定的方案、数量、路径、时刻或控制变量含义】 & 【实际单位】 \\
    【模型参数】 & 【由题目、附件、文献或前序问题确定的参数含义】 & 【实际单位】 \\
    【目标函数或评价指标】 & 【优化目标、误差指标、效益指标或约束度量的含义】 & 【单位或无量纲】 \\
    【其他高频符号】 & 【全文反复使用且无法归入上述类别的符号含义】 & 【单位】 \\
    \bottomrule
  \end{tabular}
\end{table}

\noindent\textbf{注：}填写时应保留符号的上下标、向量或矩阵格式，并确保其与模型公式、程序变量和结果表中的定义一致。
